% ==============================================================================
% ARCHITECTURE DU SYSTÈME HYBRIDE GRAPHRAG BIOMÉDICAL (LATEX / TIKZ)
% Projet PFA (5IADA - EMSI) : "Diseases Cause Oriented Analysis using RAG and Knowledge Graphs"
% Auteurs : Hamza Essamoudi • Salma Ouahman • Ghita El Hajjouji
% ==============================================================================
% Instructions d'utilisation :
% 1. En standalone : compilez directement ce fichier avec pdflatex ou xelatex.
% 2. Dans votre rapport de PFA : insérez le bloc \begin{tikzpicture}...\end{tikzpicture}
%    dans votre chapitre de conception avec % ==============================================================================
% ARCHITECTURE DU SYSTÈME HYBRIDE GRAPHRAG BIOMÉDICAL (LATEX / TIKZ)
% Projet PFA (5IADA - EMSI) : "Diseases Cause Oriented Analysis using RAG and Knowledge Graphs"
% Auteurs : Hamza Essamoudi • Salma Ouahman • Ghita El Hajjouji
% ==============================================================================
% Instructions d'utilisation :
% 1. En standalone : compilez directement ce fichier avec pdflatex ou xelatex.
% 2. Dans votre rapport de PFA : insérez le bloc \begin{tikzpicture}...\end{tikzpicture}
%    dans votre chapitre de conception avec % ==============================================================================
% ARCHITECTURE DU SYSTÈME HYBRIDE GRAPHRAG BIOMÉDICAL (LATEX / TIKZ)
% Projet PFA (5IADA - EMSI) : "Diseases Cause Oriented Analysis using RAG and Knowledge Graphs"
% Auteurs : Hamza Essamoudi • Salma Ouahman • Ghita El Hajjouji
% ==============================================================================
% Instructions d'utilisation :
% 1. En standalone : compilez directement ce fichier avec pdflatex ou xelatex.
% 2. Dans votre rapport de PFA : insérez le bloc \begin{tikzpicture}...\end{tikzpicture}
%    dans votre chapitre de conception avec % ==============================================================================
% ARCHITECTURE DU SYSTÈME HYBRIDE GRAPHRAG BIOMÉDICAL (LATEX / TIKZ)
% Projet PFA (5IADA - EMSI) : "Diseases Cause Oriented Analysis using RAG and Knowledge Graphs"
% Auteurs : Hamza Essamoudi • Salma Ouahman • Ghita El Hajjouji
% ==============================================================================
% Instructions d'utilisation :
% 1. En standalone : compilez directement ce fichier avec pdflatex ou xelatex.
% 2. Dans votre rapport de PFA : insérez le bloc \begin{tikzpicture}...\end{tikzpicture}
%    dans votre chapitre de conception avec \input{architecture_graphrag.tex}.
% ==============================================================================

\documentclass[border=15pt, 11pt]{standalone}
\usepackage[utf8]{inputenc}
\usepackage[T1]{fontenc}
\usepackage{lmodern}
\usepackage{amsmath, amssymb}
\usepackage{xcolor}
\usepackage{tikz}

% Bibliothèques TikZ nécessaires
\usetikzlibrary{shapes.geometric, arrows.meta, positioning, calc, fit, backgrounds, shadows.blur}

% Palette de couleurs professionnelles (Charte Slate Navy, Teal, Indigo & Emerald)
\definecolor{SlateNavy}{HTML}{0B1120}
\definecolor{SlateCard}{HTML}{0F172A}
\definecolor{SlateSubCard}{HTML}{1E293B}
\definecolor{TealPrimary}{HTML}{042F2E}
\definecolor{TealBorder}{HTML}{2DD4BF}
\definecolor{TealText}{HTML}{5EEAD4}
\definecolor{IndigoPrimary}{HTML}{1E1B4B}
\definecolor{IndigoBorder}{HTML}{818CF8}
\definecolor{IndigoText}{HTML}{A5B4FC}
\definecolor{SkyPrimary}{HTML}{0369A1}
\definecolor{SkyBorder}{HTML}{38BDF8}
\definecolor{AmberPrimary}{HTML}{451A03}
\definecolor{AmberBorder}{HTML}{F59E0B}
\definecolor{AmberText}{HTML}{FDE68A}
\definecolor{EmeraldPrimary}{HTML}{022C22}
\definecolor{EmeraldBorder}{HTML}{10B981}
\definecolor{EmeraldText}{HTML}{6EE7B7}
\definecolor{TextLight}{HTML}{F8FAFC}
\definecolor{TextMuted}{HTML}{94A3B8}
\definecolor{BadgeBg}{HTML}{030712}

\begin{document}

\begin{tikzpicture}[
    font=\sffamily,
    >=Stealth,
    node distance=1.2cm and 1.5cm,
    % Styles des blocs
    titleBox/.style={
        text=TextLight,
        align=center,
        font=\bfseries\Large
    },
    subtitleBox/.style={
        text=SkyBorder,
        align=center,
        font=\bfseries\normalsize
    },
    metaBox/.style={
        text=TextMuted,
        align=center,
        font=\footnotesize\itshape
    },
    stepCard/.style={
        draw=#1,
        line width=1.5pt,
        rounded corners=8pt,
        inner sep=12pt,
        fill=SlateCard,
        blur shadow={shadow blur steps=5, shadow xshift=2pt, shadow yshift=-2pt, shadow opacity=40}
    },
    innerCard/.style={
        draw=#1,
        line width=1pt,
        rounded corners=5pt,
        inner sep=8pt,
        fill=SlateSubCard,
        text=TextLight,
        align=left,
        font=\footnotesize
    },
    badge/.style={
        draw=#1,
        line width=0.8pt,
        rounded corners=3pt,
        fill=BadgeBg,
        text=#1,
        font=\bfseries\scriptsize,
        inner sep=3pt
    },
    flowArrow/.style={
        ->,
        line width=2pt,
        draw=#1,
        shorten >=3pt,
        shorten <=3pt
    },
    flowLabel/.style={
        font=\bfseries\scriptsize,
        fill=SlateNavy,
        inner sep=2pt,
        rounded corners=2pt,
        text=#1
    }
]

    % -------------------------------------------------------------------------
    % TITRE ET EN-TÊTE PRINCIPAL
    % -------------------------------------------------------------------------
    \node[titleBox] (mainTitle) at (10, 15.2) {ARCHITECTURE DU PIPELINE HYBRIDE GRAPHRAG BIOMÉDICAL};
    \node[subtitleBox, below=0.15cm of mainTitle] (subTitle) {
        Hamza Essamoudi $\bullet$ Salma Ouahman $\bullet$ Ghita El Hajjouji \quad$|$\quad Projet PFA --- 5IADA (EMSI)
    };
    \node[metaBox, below=0.12cm of subTitle] (metaLine) {
        Analyse Causale Multi-Échelles : Maladie d'Alzheimer (DOID:10652) $\bullet$ Diabète Type 2 (DOID:9352) $\bullet$ Ontologie Human Lifestyle
    };

    % -------------------------------------------------------------------------
    % ÉTAPE 1 : ACQUISITION & ETL MULTI-ONTOLOGIES
    # -------------------------------------------------------------------------
    \node[innerCard=SkyBorder, text width=6.2cm] (etl_sub1) at (3, 11.2) {
        \textbf{\textcolor{SkyBorder}{\normalsize 1.1 Réseaux Biomédicaux}}\\[3pt]
        $\bullet$ \textbf{PrimeKG-Plus} : Ontologies \& Mappings UMLS/MeSH\\
        $\bullet$ \textbf{Hetionet v1.0} : 2,25M arêtes relationnelles\\
        $\bullet$ Harmonisation : DOID, Entrez, MeSH, Reactome\\[4pt]
        \begin{tikzpicture}
            \node[badge=SkyBorder] {Python 3.12 | Pandas | Git LFS};
        \end{tikzpicture}
    };

    \node[innerCard=TealBorder, text width=6.2cm, below=0.45cm of etl_sub1] (etl_sub2) {
        \textbf{\textcolor{TealBorder}{\normalsize 1.2 Human Lifestyle Ontology (Nouveau)}}\\[3pt]
        $\bullet$ \textbf{15 N\oe uds} hiérarchiques (MeSH \texttt{D008019})\\
        $\bullet$ \textbf{40 Arêtes causales} (\texttt{PROTECTS}, \texttt{AGGRAVATES})\\
        $\bullet$ Facteurs modifiables : Diet, Exercise, Sleep, Stress\\[4pt]
        \begin{tikzpicture}
            \node[badge=TealBorder] {MeSH D008019 | UMLS C0023674 | enrich\_lifestyle.py};
        \end{tikzpicture}
    };

    \node[innerCard=SkyBorder, text width=6.2cm, below=0.45cm of etl_sub2] (etl_sub3) {
        \textbf{\textcolor{SkyBorder}{\normalsize 1.3 Littérature Médicale Certifiée}}\\[3pt]
        $\bullet$ \textbf{NIH MedQuAD} : 47 457 Q/R cliniques révisées\\
        $\bullet$ Dossiers ciblés : NINDS, GHR, NIDDK, SeniorHealth\\
        $\bullet$ Chunking contextualisé : \textbf{719 Chunks} (190 Lifestyle)\\[4pt]
        \begin{tikzpicture}
            \node[badge=SkyBorder] {XML Parsing | Chunking Contextualisé};
        \end{tikzpicture}
    };

    \begin{scope}[on background layer]
        \node[stepCard=SkyBorder, fit={(etl_sub1) (etl_sub2) (etl_sub3)}, label={[SkyBorder, font=\bfseries\normalsize]above:ÉTAPE 1 : ACQUISITION \& ETL MULTI-ONTOLOGIES}] (card_step1) {};
    \end{scope}

    % -------------------------------------------------------------------------
    % ÉTAPE 2A : KNOWLEDGE GRAPH CAUSAL (NEO4J)
    % -------------------------------------------------------------------------
    \node[innerCard=TealBorder, fill=TealPrimary, text width=6.8cm] (kg_box) at (11.5, 11.5) {
        \textbf{\textcolor{TealText}{\large 2A. Knowledge Graph Causal (Neo4j)}}\\[3pt]
        $\bullet$ \textbf{4 518 N\oe uds} uniques (9 labels hiérarchisés)\\
        $\bullet$ \textbf{19 951 Relations} causales \& fonctionnelles (17 types)\\
        $\bullet$ \textbf{Axe Diabète-Alzheimer} : 36 gènes partagés (\textit{Type 3 Diabetes})\\
        $\bullet$ Ingestion batch \texttt{UNWIND} \& Contraintes d'unicité\\
        $\bullet$ Traçabilité $k$-hop : Lifestyle $\to$ Gènes $\to$ Voies $\to$ Maladie\\[4pt]
        \begin{tikzpicture}
            \node[badge=TealBorder] {Neo4j 5.26 | Cypher | APOC | Docker Compose};
        \end{tikzpicture}
    };

    % -------------------------------------------------------------------------
    % ÉTAPE 2B : BASE VECTORIELLE CLINIQUE DENSE
    % -------------------------------------------------------------------------
    \node[innerCard=IndigoBorder, fill=IndigoPrimary, text width=6.8cm, below=0.6cm of kg_box] (vec_box) {
        \textbf{\textcolor{IndigoText}{\large 2B. Base Vectorielle Clinique Dense}}\\[3pt]
        $\bullet$ \textbf{719 Chunks} contextuels MedQuAD indexés\\
        $\bullet$ Modèle biomédical \textbf{PubMedBERT} (768 dimensions)\\
        $\bullet$ Métadonnées cliniques : \texttt{disease}, \texttt{qtype}, \texttt{focus}\\
        $\bullet$ \textbf{Dual Vector Store} :\\
        \quad - \textbf{ChromaDB} : Moteur embarqué local sans dépendance\\
        \quad - \textbf{Qdrant} : Microservice Docker pour la production\\[4pt]
        \begin{tikzpicture}
            \node[badge=IndigoBorder] {PubMedBERT (768d) | ChromaDB | Qdrant Docker};
        \end{tikzpicture}
    };

    \begin{scope}[on background layer]
        \node[stepCard=TealBorder, fit={(kg_box) (vec_box)}, label={[TealBorder, font=\bfseries\normalsize]above:ÉTAPE 2 : STOCKAGE \& INDEXATION DUAL-CORE}] (card_step2) {};
    \end{scope}

    % -------------------------------------------------------------------------
    % ÉTAPE 3 : RETRIEVER HYBRIDE & FUSION DE CONTEXTE
    % -------------------------------------------------------------------------
    \node[innerCard=AmberBorder, fill=AmberPrimary, text width=6.6cm] (rag_box1) at (19.8, 11.5) {
        \textbf{\textcolor{AmberText}{\large 3.1 Retrieval Conjoint}}\\[3pt]
        $\bullet$ \textbf{BiomedicalEntityLinker} : 30+ alias lifestyle \& gènes\\
        $\bullet$ \textbf{Graphe $k$-Hop} : (Lifestyle) $\to$ (BDNF/TNF) $\to$ (Voies)\\
        $\bullet$ \textbf{Axe Diabète-Alzheimer} : 36 gènes communs (IRS1, AKT1)\\
        $\bullet$ \textbf{Recherche Vectorielle} : Top-$k$ chunks certifiés NIH\\
        $\bullet$ Score combiné normalisé de pertinence\\[4pt]
        \begin{tikzpicture}
            \node[badge=AmberBorder] {HybridRetriever | LocalGraphEngine};
        \end{tikzpicture}
    };

    \node[innerCard=AmberBorder, fill=AmberPrimary, text width=6.6cm, below=0.6cm of rag_box1] (rag_box2) {
        \textbf{\textcolor{AmberText}{\large 3.2 Fusion de Contexte}}\\[3pt]
        $\bullet$ Linéarisation ordonnée des triplets \& chaînes causales\\
        $\bullet$ Fusion conjointe : Graphe formel + Extraits certifiés NIH\\
        $\bullet$ Élimination documentée des hallucinations médicales\\
        $\bullet$ Éradication des ruptures déductives et faux liens\\[4pt]
        \begin{tikzpicture}
            \node[badge=AmberBorder] {ContextFusionEngine | Pydantic Schema};
        \end{tikzpicture}
    };

    \begin{scope}[on background layer]
        \node[stepCard=AmberBorder, fit={(rag_box1) (rag_box2)}, label={[AmberBorder, font=\bfseries\normalsize]above:ÉTAPE 3 : MOTEUR HYBRIDE GRAPHRAG}] (card_step3) {};
    \end{scope}

    % -------------------------------------------------------------------------
    % ÉTAPES 4 & 5 : RAISONNEMENT CAUSAL, BENCHMARK & APPLICATION STREAMLIT
    % -------------------------------------------------------------------------
    \node[innerCard=EmeraldBorder, fill=EmeraldPrimary, text width=6.2cm] (llm_box1) at (3.2, 1.8) {
        \textbf{\textcolor{EmeraldText}{\large 4. Moteur de Synthèse Causale}}\\[3pt]
        1. Déclencheurs étiologiques (APP, PSEN1, APOE)\\
        2. Cascades intermédiaires (Résistance insuline, TNF)\\
        3. Impact Lifestyle : Neuroprotection (BDNF)\\
        4. Phénotypes \& déclin cognitif (Démence, mémoire)\\
        5. Repositionnement (Donepezil, Memantine, Diète)\\
        6. Traçabilité absolue : Triplets graphe \& DocIDs NIH\\[4pt]
        \begin{tikzpicture}
            \node[badge=EmeraldBorder] {CausalGenerator | LLM Routing};
        \end{tikzpicture}
    };

    \node[innerCard=EmeraldBorder, fill=EmeraldPrimary, text width=6.6cm] (llm_box2) at (11.5, 1.8) {
        \textbf{\textcolor{EmeraldText}{\large 5. Benchmark d'Évaluation}}\\[3pt]
        $\bullet$ \textbf{Chaînes causales tracées} : 0 (Std) vs \textbf{23} (GraphRAG)\\
        $\bullet$ \textbf{Couverture des voies métaboliques} : \textbf{+100\%}\\
        $\bullet$ \textbf{Triplets \& Citations d'évidences} : 15 vs \textbf{35}\\
        $\bullet$ Connexions Diabète-Alzheimer : 0 vs 100\% identifiées\\
        $\bullet$ Fichier certifié : \texttt{data/processed/evaluation\_results.json}\\[4pt]
        \begin{tikzpicture}
            \node[badge=EmeraldBorder] {evaluate\_rag.py | Benchmark Suite};
        \end{tikzpicture}
    };

    \node[innerCard=EmeraldBorder, fill=EmeraldPrimary, text width=6.2cm] (llm_box3) at (19.8, 1.8) {
        \textbf{\textcolor{EmeraldText}{\large Interfaces \& Déploiement}}\\[3pt]
        $\bullet$ \textbf{Application Web Streamlit} (\texttt{app.py} \& \texttt{run\_app.bat}) :\\
        \quad - Chatbot GraphRAG avec streaming \& citations\\
        \quad - Visualiseur (36 gènes partagés \& mode de vie)\\
        \quad - Dashboard comparatif Standard RAG vs GraphRAG\\
        $\bullet$ \textbf{CLI Interactif} (\texttt{src/rag/cli.py}) \& Démo (\texttt{demo.py})\\
        $\bullet$ Dépôt GitHub : \texttt{biomedical-graphrag}\\[4pt]
        \begin{tikzpicture}
            \node[badge=EmeraldBorder] {Streamlit (app.py) | CLI | GitHub CI};
        \end{tikzpicture}
    };

    \begin{scope}[on background layer]
        \node[stepCard=EmeraldBorder, fit={(llm_box1) (llm_box2) (llm_box3)}, label={[EmeraldBorder, font=\bfseries\normalsize]above:ÉTAPE 4 \& 5 : RAISONNEMENT CAUSAL MULTI-ÉCHELLES, BENCHMARK \& DÉPLOIEMENT STREAMLIT}] (card_step4) {};
    \end{scope}

    % -------------------------------------------------------------------------
    % FLÈCHES DE FLUX DIRECTIONNELLES (DATAFLOW)
    % -------------------------------------------------------------------------
    % ETL vers Graphe
    \draw[flowArrow=TealBorder] (etl_sub1.east) -- (kg_box.west)
        node[midway, flowLabel=TealBorder] {Triplets CSV};

    % ETL vers Vector Store
    \draw[flowArrow=IndigoBorder] (etl_sub3.east) -- (vec_box.west)
        node[midway, flowLabel=IndigoBorder] {Chunks XML};

    % Graphe vers Hybrid Retriever
    \draw[flowArrow=TealBorder] (kg_box.east) -- (rag_box1.west)
        node[midway, flowLabel=TealBorder] {Chemins $k$-Hop \& 36 Gènes};

    % Vector Store vers Hybrid Retriever
    \draw[flowArrow=IndigoBorder] (vec_box.east) -- (rag_box2.west)
        node[midway, flowLabel=IndigoBorder] {Top-$k$ Chunks (PubMedBERT)};

    % Hybrid Retriever vers Causal Generator
    \draw[flowArrow=AmberBorder] (card_step3.south) -- ++(0,-0.4) -| (llm_box3.north)
        node[pos=0.25, flowLabel=AmberBorder] {Contexte Causal Fusionné};

\end{tikzpicture}

\end{document}
.
% ==============================================================================

\documentclass[border=15pt, 11pt]{standalone}
\usepackage[utf8]{inputenc}
\usepackage[T1]{fontenc}
\usepackage{lmodern}
\usepackage{amsmath, amssymb}
\usepackage{xcolor}
\usepackage{tikz}

% Bibliothèques TikZ nécessaires
\usetikzlibrary{shapes.geometric, arrows.meta, positioning, calc, fit, backgrounds, shadows.blur}

% Palette de couleurs professionnelles (Charte Slate Navy, Teal, Indigo & Emerald)
\definecolor{SlateNavy}{HTML}{0B1120}
\definecolor{SlateCard}{HTML}{0F172A}
\definecolor{SlateSubCard}{HTML}{1E293B}
\definecolor{TealPrimary}{HTML}{042F2E}
\definecolor{TealBorder}{HTML}{2DD4BF}
\definecolor{TealText}{HTML}{5EEAD4}
\definecolor{IndigoPrimary}{HTML}{1E1B4B}
\definecolor{IndigoBorder}{HTML}{818CF8}
\definecolor{IndigoText}{HTML}{A5B4FC}
\definecolor{SkyPrimary}{HTML}{0369A1}
\definecolor{SkyBorder}{HTML}{38BDF8}
\definecolor{AmberPrimary}{HTML}{451A03}
\definecolor{AmberBorder}{HTML}{F59E0B}
\definecolor{AmberText}{HTML}{FDE68A}
\definecolor{EmeraldPrimary}{HTML}{022C22}
\definecolor{EmeraldBorder}{HTML}{10B981}
\definecolor{EmeraldText}{HTML}{6EE7B7}
\definecolor{TextLight}{HTML}{F8FAFC}
\definecolor{TextMuted}{HTML}{94A3B8}
\definecolor{BadgeBg}{HTML}{030712}

\begin{document}

\begin{tikzpicture}[
    font=\sffamily,
    >=Stealth,
    node distance=1.2cm and 1.5cm,
    % Styles des blocs
    titleBox/.style={
        text=TextLight,
        align=center,
        font=\bfseries\Large
    },
    subtitleBox/.style={
        text=SkyBorder,
        align=center,
        font=\bfseries\normalsize
    },
    metaBox/.style={
        text=TextMuted,
        align=center,
        font=\footnotesize\itshape
    },
    stepCard/.style={
        draw=#1,
        line width=1.5pt,
        rounded corners=8pt,
        inner sep=12pt,
        fill=SlateCard,
        blur shadow={shadow blur steps=5, shadow xshift=2pt, shadow yshift=-2pt, shadow opacity=40}
    },
    innerCard/.style={
        draw=#1,
        line width=1pt,
        rounded corners=5pt,
        inner sep=8pt,
        fill=SlateSubCard,
        text=TextLight,
        align=left,
        font=\footnotesize
    },
    badge/.style={
        draw=#1,
        line width=0.8pt,
        rounded corners=3pt,
        fill=BadgeBg,
        text=#1,
        font=\bfseries\scriptsize,
        inner sep=3pt
    },
    flowArrow/.style={
        ->,
        line width=2pt,
        draw=#1,
        shorten >=3pt,
        shorten <=3pt
    },
    flowLabel/.style={
        font=\bfseries\scriptsize,
        fill=SlateNavy,
        inner sep=2pt,
        rounded corners=2pt,
        text=#1
    }
]

    % -------------------------------------------------------------------------
    % TITRE ET EN-TÊTE PRINCIPAL
    % -------------------------------------------------------------------------
    \node[titleBox] (mainTitle) at (10, 15.2) {ARCHITECTURE DU PIPELINE HYBRIDE GRAPHRAG BIOMÉDICAL};
    \node[subtitleBox, below=0.15cm of mainTitle] (subTitle) {
        Hamza Essamoudi $\bullet$ Salma Ouahman $\bullet$ Ghita El Hajjouji \quad$|$\quad Projet PFA --- 5IADA (EMSI)
    };
    \node[metaBox, below=0.12cm of subTitle] (metaLine) {
        Analyse Causale Multi-Échelles : Maladie d'Alzheimer (DOID:10652) $\bullet$ Diabète Type 2 (DOID:9352) $\bullet$ Ontologie Human Lifestyle
    };

    % -------------------------------------------------------------------------
    % ÉTAPE 1 : ACQUISITION & ETL MULTI-ONTOLOGIES
    # -------------------------------------------------------------------------
    \node[innerCard=SkyBorder, text width=6.2cm] (etl_sub1) at (3, 11.2) {
        \textbf{\textcolor{SkyBorder}{\normalsize 1.1 Réseaux Biomédicaux}}\\[3pt]
        $\bullet$ \textbf{PrimeKG-Plus} : Ontologies \& Mappings UMLS/MeSH\\
        $\bullet$ \textbf{Hetionet v1.0} : 2,25M arêtes relationnelles\\
        $\bullet$ Harmonisation : DOID, Entrez, MeSH, Reactome\\[4pt]
        \begin{tikzpicture}
            \node[badge=SkyBorder] {Python 3.12 | Pandas | Git LFS};
        \end{tikzpicture}
    };

    \node[innerCard=TealBorder, text width=6.2cm, below=0.45cm of etl_sub1] (etl_sub2) {
        \textbf{\textcolor{TealBorder}{\normalsize 1.2 Human Lifestyle Ontology (Nouveau)}}\\[3pt]
        $\bullet$ \textbf{15 N\oe uds} hiérarchiques (MeSH \texttt{D008019})\\
        $\bullet$ \textbf{40 Arêtes causales} (\texttt{PROTECTS}, \texttt{AGGRAVATES})\\
        $\bullet$ Facteurs modifiables : Diet, Exercise, Sleep, Stress\\[4pt]
        \begin{tikzpicture}
            \node[badge=TealBorder] {MeSH D008019 | UMLS C0023674 | enrich\_lifestyle.py};
        \end{tikzpicture}
    };

    \node[innerCard=SkyBorder, text width=6.2cm, below=0.45cm of etl_sub2] (etl_sub3) {
        \textbf{\textcolor{SkyBorder}{\normalsize 1.3 Littérature Médicale Certifiée}}\\[3pt]
        $\bullet$ \textbf{NIH MedQuAD} : 47 457 Q/R cliniques révisées\\
        $\bullet$ Dossiers ciblés : NINDS, GHR, NIDDK, SeniorHealth\\
        $\bullet$ Chunking contextualisé : \textbf{719 Chunks} (190 Lifestyle)\\[4pt]
        \begin{tikzpicture}
            \node[badge=SkyBorder] {XML Parsing | Chunking Contextualisé};
        \end{tikzpicture}
    };

    \begin{scope}[on background layer]
        \node[stepCard=SkyBorder, fit={(etl_sub1) (etl_sub2) (etl_sub3)}, label={[SkyBorder, font=\bfseries\normalsize]above:ÉTAPE 1 : ACQUISITION \& ETL MULTI-ONTOLOGIES}] (card_step1) {};
    \end{scope}

    % -------------------------------------------------------------------------
    % ÉTAPE 2A : KNOWLEDGE GRAPH CAUSAL (NEO4J)
    % -------------------------------------------------------------------------
    \node[innerCard=TealBorder, fill=TealPrimary, text width=6.8cm] (kg_box) at (11.5, 11.5) {
        \textbf{\textcolor{TealText}{\large 2A. Knowledge Graph Causal (Neo4j)}}\\[3pt]
        $\bullet$ \textbf{4 518 N\oe uds} uniques (9 labels hiérarchisés)\\
        $\bullet$ \textbf{19 951 Relations} causales \& fonctionnelles (17 types)\\
        $\bullet$ \textbf{Axe Diabète-Alzheimer} : 36 gènes partagés (\textit{Type 3 Diabetes})\\
        $\bullet$ Ingestion batch \texttt{UNWIND} \& Contraintes d'unicité\\
        $\bullet$ Traçabilité $k$-hop : Lifestyle $\to$ Gènes $\to$ Voies $\to$ Maladie\\[4pt]
        \begin{tikzpicture}
            \node[badge=TealBorder] {Neo4j 5.26 | Cypher | APOC | Docker Compose};
        \end{tikzpicture}
    };

    % -------------------------------------------------------------------------
    % ÉTAPE 2B : BASE VECTORIELLE CLINIQUE DENSE
    % -------------------------------------------------------------------------
    \node[innerCard=IndigoBorder, fill=IndigoPrimary, text width=6.8cm, below=0.6cm of kg_box] (vec_box) {
        \textbf{\textcolor{IndigoText}{\large 2B. Base Vectorielle Clinique Dense}}\\[3pt]
        $\bullet$ \textbf{719 Chunks} contextuels MedQuAD indexés\\
        $\bullet$ Modèle biomédical \textbf{PubMedBERT} (768 dimensions)\\
        $\bullet$ Métadonnées cliniques : \texttt{disease}, \texttt{qtype}, \texttt{focus}\\
        $\bullet$ \textbf{Dual Vector Store} :\\
        \quad - \textbf{ChromaDB} : Moteur embarqué local sans dépendance\\
        \quad - \textbf{Qdrant} : Microservice Docker pour la production\\[4pt]
        \begin{tikzpicture}
            \node[badge=IndigoBorder] {PubMedBERT (768d) | ChromaDB | Qdrant Docker};
        \end{tikzpicture}
    };

    \begin{scope}[on background layer]
        \node[stepCard=TealBorder, fit={(kg_box) (vec_box)}, label={[TealBorder, font=\bfseries\normalsize]above:ÉTAPE 2 : STOCKAGE \& INDEXATION DUAL-CORE}] (card_step2) {};
    \end{scope}

    % -------------------------------------------------------------------------
    % ÉTAPE 3 : RETRIEVER HYBRIDE & FUSION DE CONTEXTE
    % -------------------------------------------------------------------------
    \node[innerCard=AmberBorder, fill=AmberPrimary, text width=6.6cm] (rag_box1) at (19.8, 11.5) {
        \textbf{\textcolor{AmberText}{\large 3.1 Retrieval Conjoint}}\\[3pt]
        $\bullet$ \textbf{BiomedicalEntityLinker} : 30+ alias lifestyle \& gènes\\
        $\bullet$ \textbf{Graphe $k$-Hop} : (Lifestyle) $\to$ (BDNF/TNF) $\to$ (Voies)\\
        $\bullet$ \textbf{Axe Diabète-Alzheimer} : 36 gènes communs (IRS1, AKT1)\\
        $\bullet$ \textbf{Recherche Vectorielle} : Top-$k$ chunks certifiés NIH\\
        $\bullet$ Score combiné normalisé de pertinence\\[4pt]
        \begin{tikzpicture}
            \node[badge=AmberBorder] {HybridRetriever | LocalGraphEngine};
        \end{tikzpicture}
    };

    \node[innerCard=AmberBorder, fill=AmberPrimary, text width=6.6cm, below=0.6cm of rag_box1] (rag_box2) {
        \textbf{\textcolor{AmberText}{\large 3.2 Fusion de Contexte}}\\[3pt]
        $\bullet$ Linéarisation ordonnée des triplets \& chaînes causales\\
        $\bullet$ Fusion conjointe : Graphe formel + Extraits certifiés NIH\\
        $\bullet$ Élimination documentée des hallucinations médicales\\
        $\bullet$ Éradication des ruptures déductives et faux liens\\[4pt]
        \begin{tikzpicture}
            \node[badge=AmberBorder] {ContextFusionEngine | Pydantic Schema};
        \end{tikzpicture}
    };

    \begin{scope}[on background layer]
        \node[stepCard=AmberBorder, fit={(rag_box1) (rag_box2)}, label={[AmberBorder, font=\bfseries\normalsize]above:ÉTAPE 3 : MOTEUR HYBRIDE GRAPHRAG}] (card_step3) {};
    \end{scope}

    % -------------------------------------------------------------------------
    % ÉTAPES 4 & 5 : RAISONNEMENT CAUSAL, BENCHMARK & APPLICATION STREAMLIT
    % -------------------------------------------------------------------------
    \node[innerCard=EmeraldBorder, fill=EmeraldPrimary, text width=6.2cm] (llm_box1) at (3.2, 1.8) {
        \textbf{\textcolor{EmeraldText}{\large 4. Moteur de Synthèse Causale}}\\[3pt]
        1. Déclencheurs étiologiques (APP, PSEN1, APOE)\\
        2. Cascades intermédiaires (Résistance insuline, TNF)\\
        3. Impact Lifestyle : Neuroprotection (BDNF)\\
        4. Phénotypes \& déclin cognitif (Démence, mémoire)\\
        5. Repositionnement (Donepezil, Memantine, Diète)\\
        6. Traçabilité absolue : Triplets graphe \& DocIDs NIH\\[4pt]
        \begin{tikzpicture}
            \node[badge=EmeraldBorder] {CausalGenerator | LLM Routing};
        \end{tikzpicture}
    };

    \node[innerCard=EmeraldBorder, fill=EmeraldPrimary, text width=6.6cm] (llm_box2) at (11.5, 1.8) {
        \textbf{\textcolor{EmeraldText}{\large 5. Benchmark d'Évaluation}}\\[3pt]
        $\bullet$ \textbf{Chaînes causales tracées} : 0 (Std) vs \textbf{23} (GraphRAG)\\
        $\bullet$ \textbf{Couverture des voies métaboliques} : \textbf{+100\%}\\
        $\bullet$ \textbf{Triplets \& Citations d'évidences} : 15 vs \textbf{35}\\
        $\bullet$ Connexions Diabète-Alzheimer : 0 vs 100\% identifiées\\
        $\bullet$ Fichier certifié : \texttt{data/processed/evaluation\_results.json}\\[4pt]
        \begin{tikzpicture}
            \node[badge=EmeraldBorder] {evaluate\_rag.py | Benchmark Suite};
        \end{tikzpicture}
    };

    \node[innerCard=EmeraldBorder, fill=EmeraldPrimary, text width=6.2cm] (llm_box3) at (19.8, 1.8) {
        \textbf{\textcolor{EmeraldText}{\large Interfaces \& Déploiement}}\\[3pt]
        $\bullet$ \textbf{Application Web Streamlit} (\texttt{app.py} \& \texttt{run\_app.bat}) :\\
        \quad - Chatbot GraphRAG avec streaming \& citations\\
        \quad - Visualiseur (36 gènes partagés \& mode de vie)\\
        \quad - Dashboard comparatif Standard RAG vs GraphRAG\\
        $\bullet$ \textbf{CLI Interactif} (\texttt{src/rag/cli.py}) \& Démo (\texttt{demo.py})\\
        $\bullet$ Dépôt GitHub : \texttt{biomedical-graphrag}\\[4pt]
        \begin{tikzpicture}
            \node[badge=EmeraldBorder] {Streamlit (app.py) | CLI | GitHub CI};
        \end{tikzpicture}
    };

    \begin{scope}[on background layer]
        \node[stepCard=EmeraldBorder, fit={(llm_box1) (llm_box2) (llm_box3)}, label={[EmeraldBorder, font=\bfseries\normalsize]above:ÉTAPE 4 \& 5 : RAISONNEMENT CAUSAL MULTI-ÉCHELLES, BENCHMARK \& DÉPLOIEMENT STREAMLIT}] (card_step4) {};
    \end{scope}

    % -------------------------------------------------------------------------
    % FLÈCHES DE FLUX DIRECTIONNELLES (DATAFLOW)
    % -------------------------------------------------------------------------
    % ETL vers Graphe
    \draw[flowArrow=TealBorder] (etl_sub1.east) -- (kg_box.west)
        node[midway, flowLabel=TealBorder] {Triplets CSV};

    % ETL vers Vector Store
    \draw[flowArrow=IndigoBorder] (etl_sub3.east) -- (vec_box.west)
        node[midway, flowLabel=IndigoBorder] {Chunks XML};

    % Graphe vers Hybrid Retriever
    \draw[flowArrow=TealBorder] (kg_box.east) -- (rag_box1.west)
        node[midway, flowLabel=TealBorder] {Chemins $k$-Hop \& 36 Gènes};

    % Vector Store vers Hybrid Retriever
    \draw[flowArrow=IndigoBorder] (vec_box.east) -- (rag_box2.west)
        node[midway, flowLabel=IndigoBorder] {Top-$k$ Chunks (PubMedBERT)};

    % Hybrid Retriever vers Causal Generator
    \draw[flowArrow=AmberBorder] (card_step3.south) -- ++(0,-0.4) -| (llm_box3.north)
        node[pos=0.25, flowLabel=AmberBorder] {Contexte Causal Fusionné};

\end{tikzpicture}

\end{document}
.
% ==============================================================================

\documentclass[border=15pt, 11pt]{standalone}
\usepackage[utf8]{inputenc}
\usepackage[T1]{fontenc}
\usepackage{lmodern}
\usepackage{amsmath, amssymb}
\usepackage{xcolor}
\usepackage{tikz}

% Bibliothèques TikZ nécessaires
\usetikzlibrary{shapes.geometric, arrows.meta, positioning, calc, fit, backgrounds, shadows.blur}

% Palette de couleurs professionnelles (Charte Slate Navy, Teal, Indigo & Emerald)
\definecolor{SlateNavy}{HTML}{0B1120}
\definecolor{SlateCard}{HTML}{0F172A}
\definecolor{SlateSubCard}{HTML}{1E293B}
\definecolor{TealPrimary}{HTML}{042F2E}
\definecolor{TealBorder}{HTML}{2DD4BF}
\definecolor{TealText}{HTML}{5EEAD4}
\definecolor{IndigoPrimary}{HTML}{1E1B4B}
\definecolor{IndigoBorder}{HTML}{818CF8}
\definecolor{IndigoText}{HTML}{A5B4FC}
\definecolor{SkyPrimary}{HTML}{0369A1}
\definecolor{SkyBorder}{HTML}{38BDF8}
\definecolor{AmberPrimary}{HTML}{451A03}
\definecolor{AmberBorder}{HTML}{F59E0B}
\definecolor{AmberText}{HTML}{FDE68A}
\definecolor{EmeraldPrimary}{HTML}{022C22}
\definecolor{EmeraldBorder}{HTML}{10B981}
\definecolor{EmeraldText}{HTML}{6EE7B7}
\definecolor{TextLight}{HTML}{F8FAFC}
\definecolor{TextMuted}{HTML}{94A3B8}
\definecolor{BadgeBg}{HTML}{030712}

\begin{document}

\begin{tikzpicture}[
    font=\sffamily,
    >=Stealth,
    node distance=1.2cm and 1.5cm,
    % Styles des blocs
    titleBox/.style={
        text=TextLight,
        align=center,
        font=\bfseries\Large
    },
    subtitleBox/.style={
        text=SkyBorder,
        align=center,
        font=\bfseries\normalsize
    },
    metaBox/.style={
        text=TextMuted,
        align=center,
        font=\footnotesize\itshape
    },
    stepCard/.style={
        draw=#1,
        line width=1.5pt,
        rounded corners=8pt,
        inner sep=12pt,
        fill=SlateCard,
        blur shadow={shadow blur steps=5, shadow xshift=2pt, shadow yshift=-2pt, shadow opacity=40}
    },
    innerCard/.style={
        draw=#1,
        line width=1pt,
        rounded corners=5pt,
        inner sep=8pt,
        fill=SlateSubCard,
        text=TextLight,
        align=left,
        font=\footnotesize
    },
    badge/.style={
        draw=#1,
        line width=0.8pt,
        rounded corners=3pt,
        fill=BadgeBg,
        text=#1,
        font=\bfseries\scriptsize,
        inner sep=3pt
    },
    flowArrow/.style={
        ->,
        line width=2pt,
        draw=#1,
        shorten >=3pt,
        shorten <=3pt
    },
    flowLabel/.style={
        font=\bfseries\scriptsize,
        fill=SlateNavy,
        inner sep=2pt,
        rounded corners=2pt,
        text=#1
    }
]

    % -------------------------------------------------------------------------
    % TITRE ET EN-TÊTE PRINCIPAL
    % -------------------------------------------------------------------------
    \node[titleBox] (mainTitle) at (10, 15.2) {ARCHITECTURE DU PIPELINE HYBRIDE GRAPHRAG BIOMÉDICAL};
    \node[subtitleBox, below=0.15cm of mainTitle] (subTitle) {
        Hamza Essamoudi $\bullet$ Salma Ouahman $\bullet$ Ghita El Hajjouji \quad$|$\quad Projet PFA --- 5IADA (EMSI)
    };
    \node[metaBox, below=0.12cm of subTitle] (metaLine) {
        Analyse Causale Multi-Échelles : Maladie d'Alzheimer (DOID:10652) $\bullet$ Diabète Type 2 (DOID:9352) $\bullet$ Ontologie Human Lifestyle
    };

    % -------------------------------------------------------------------------
    % ÉTAPE 1 : ACQUISITION & ETL MULTI-ONTOLOGIES
    # -------------------------------------------------------------------------
    \node[innerCard=SkyBorder, text width=6.2cm] (etl_sub1) at (3, 11.2) {
        \textbf{\textcolor{SkyBorder}{\normalsize 1.1 Réseaux Biomédicaux}}\\[3pt]
        $\bullet$ \textbf{PrimeKG-Plus} : Ontologies \& Mappings UMLS/MeSH\\
        $\bullet$ \textbf{Hetionet v1.0} : 2,25M arêtes relationnelles\\
        $\bullet$ Harmonisation : DOID, Entrez, MeSH, Reactome\\[4pt]
        \begin{tikzpicture}
            \node[badge=SkyBorder] {Python 3.12 | Pandas | Git LFS};
        \end{tikzpicture}
    };

    \node[innerCard=TealBorder, text width=6.2cm, below=0.45cm of etl_sub1] (etl_sub2) {
        \textbf{\textcolor{TealBorder}{\normalsize 1.2 Human Lifestyle Ontology (Nouveau)}}\\[3pt]
        $\bullet$ \textbf{15 N\oe uds} hiérarchiques (MeSH \texttt{D008019})\\
        $\bullet$ \textbf{40 Arêtes causales} (\texttt{PROTECTS}, \texttt{AGGRAVATES})\\
        $\bullet$ Facteurs modifiables : Diet, Exercise, Sleep, Stress\\[4pt]
        \begin{tikzpicture}
            \node[badge=TealBorder] {MeSH D008019 | UMLS C0023674 | enrich\_lifestyle.py};
        \end{tikzpicture}
    };

    \node[innerCard=SkyBorder, text width=6.2cm, below=0.45cm of etl_sub2] (etl_sub3) {
        \textbf{\textcolor{SkyBorder}{\normalsize 1.3 Littérature Médicale Certifiée}}\\[3pt]
        $\bullet$ \textbf{NIH MedQuAD} : 47 457 Q/R cliniques révisées\\
        $\bullet$ Dossiers ciblés : NINDS, GHR, NIDDK, SeniorHealth\\
        $\bullet$ Chunking contextualisé : \textbf{719 Chunks} (190 Lifestyle)\\[4pt]
        \begin{tikzpicture}
            \node[badge=SkyBorder] {XML Parsing | Chunking Contextualisé};
        \end{tikzpicture}
    };

    \begin{scope}[on background layer]
        \node[stepCard=SkyBorder, fit={(etl_sub1) (etl_sub2) (etl_sub3)}, label={[SkyBorder, font=\bfseries\normalsize]above:ÉTAPE 1 : ACQUISITION \& ETL MULTI-ONTOLOGIES}] (card_step1) {};
    \end{scope}

    % -------------------------------------------------------------------------
    % ÉTAPE 2A : KNOWLEDGE GRAPH CAUSAL (NEO4J)
    % -------------------------------------------------------------------------
    \node[innerCard=TealBorder, fill=TealPrimary, text width=6.8cm] (kg_box) at (11.5, 11.5) {
        \textbf{\textcolor{TealText}{\large 2A. Knowledge Graph Causal (Neo4j)}}\\[3pt]
        $\bullet$ \textbf{4 518 N\oe uds} uniques (9 labels hiérarchisés)\\
        $\bullet$ \textbf{19 951 Relations} causales \& fonctionnelles (17 types)\\
        $\bullet$ \textbf{Axe Diabète-Alzheimer} : 36 gènes partagés (\textit{Type 3 Diabetes})\\
        $\bullet$ Ingestion batch \texttt{UNWIND} \& Contraintes d'unicité\\
        $\bullet$ Traçabilité $k$-hop : Lifestyle $\to$ Gènes $\to$ Voies $\to$ Maladie\\[4pt]
        \begin{tikzpicture}
            \node[badge=TealBorder] {Neo4j 5.26 | Cypher | APOC | Docker Compose};
        \end{tikzpicture}
    };

    % -------------------------------------------------------------------------
    % ÉTAPE 2B : BASE VECTORIELLE CLINIQUE DENSE
    % -------------------------------------------------------------------------
    \node[innerCard=IndigoBorder, fill=IndigoPrimary, text width=6.8cm, below=0.6cm of kg_box] (vec_box) {
        \textbf{\textcolor{IndigoText}{\large 2B. Base Vectorielle Clinique Dense}}\\[3pt]
        $\bullet$ \textbf{719 Chunks} contextuels MedQuAD indexés\\
        $\bullet$ Modèle biomédical \textbf{PubMedBERT} (768 dimensions)\\
        $\bullet$ Métadonnées cliniques : \texttt{disease}, \texttt{qtype}, \texttt{focus}\\
        $\bullet$ \textbf{Dual Vector Store} :\\
        \quad - \textbf{ChromaDB} : Moteur embarqué local sans dépendance\\
        \quad - \textbf{Qdrant} : Microservice Docker pour la production\\[4pt]
        \begin{tikzpicture}
            \node[badge=IndigoBorder] {PubMedBERT (768d) | ChromaDB | Qdrant Docker};
        \end{tikzpicture}
    };

    \begin{scope}[on background layer]
        \node[stepCard=TealBorder, fit={(kg_box) (vec_box)}, label={[TealBorder, font=\bfseries\normalsize]above:ÉTAPE 2 : STOCKAGE \& INDEXATION DUAL-CORE}] (card_step2) {};
    \end{scope}

    % -------------------------------------------------------------------------
    % ÉTAPE 3 : RETRIEVER HYBRIDE & FUSION DE CONTEXTE
    % -------------------------------------------------------------------------
    \node[innerCard=AmberBorder, fill=AmberPrimary, text width=6.6cm] (rag_box1) at (19.8, 11.5) {
        \textbf{\textcolor{AmberText}{\large 3.1 Retrieval Conjoint}}\\[3pt]
        $\bullet$ \textbf{BiomedicalEntityLinker} : 30+ alias lifestyle \& gènes\\
        $\bullet$ \textbf{Graphe $k$-Hop} : (Lifestyle) $\to$ (BDNF/TNF) $\to$ (Voies)\\
        $\bullet$ \textbf{Axe Diabète-Alzheimer} : 36 gènes communs (IRS1, AKT1)\\
        $\bullet$ \textbf{Recherche Vectorielle} : Top-$k$ chunks certifiés NIH\\
        $\bullet$ Score combiné normalisé de pertinence\\[4pt]
        \begin{tikzpicture}
            \node[badge=AmberBorder] {HybridRetriever | LocalGraphEngine};
        \end{tikzpicture}
    };

    \node[innerCard=AmberBorder, fill=AmberPrimary, text width=6.6cm, below=0.6cm of rag_box1] (rag_box2) {
        \textbf{\textcolor{AmberText}{\large 3.2 Fusion de Contexte}}\\[3pt]
        $\bullet$ Linéarisation ordonnée des triplets \& chaînes causales\\
        $\bullet$ Fusion conjointe : Graphe formel + Extraits certifiés NIH\\
        $\bullet$ Élimination documentée des hallucinations médicales\\
        $\bullet$ Éradication des ruptures déductives et faux liens\\[4pt]
        \begin{tikzpicture}
            \node[badge=AmberBorder] {ContextFusionEngine | Pydantic Schema};
        \end{tikzpicture}
    };

    \begin{scope}[on background layer]
        \node[stepCard=AmberBorder, fit={(rag_box1) (rag_box2)}, label={[AmberBorder, font=\bfseries\normalsize]above:ÉTAPE 3 : MOTEUR HYBRIDE GRAPHRAG}] (card_step3) {};
    \end{scope}

    % -------------------------------------------------------------------------
    % ÉTAPES 4 & 5 : RAISONNEMENT CAUSAL, BENCHMARK & APPLICATION STREAMLIT
    % -------------------------------------------------------------------------
    \node[innerCard=EmeraldBorder, fill=EmeraldPrimary, text width=6.2cm] (llm_box1) at (3.2, 1.8) {
        \textbf{\textcolor{EmeraldText}{\large 4. Moteur de Synthèse Causale}}\\[3pt]
        1. Déclencheurs étiologiques (APP, PSEN1, APOE)\\
        2. Cascades intermédiaires (Résistance insuline, TNF)\\
        3. Impact Lifestyle : Neuroprotection (BDNF)\\
        4. Phénotypes \& déclin cognitif (Démence, mémoire)\\
        5. Repositionnement (Donepezil, Memantine, Diète)\\
        6. Traçabilité absolue : Triplets graphe \& DocIDs NIH\\[4pt]
        \begin{tikzpicture}
            \node[badge=EmeraldBorder] {CausalGenerator | LLM Routing};
        \end{tikzpicture}
    };

    \node[innerCard=EmeraldBorder, fill=EmeraldPrimary, text width=6.6cm] (llm_box2) at (11.5, 1.8) {
        \textbf{\textcolor{EmeraldText}{\large 5. Benchmark d'Évaluation}}\\[3pt]
        $\bullet$ \textbf{Chaînes causales tracées} : 0 (Std) vs \textbf{23} (GraphRAG)\\
        $\bullet$ \textbf{Couverture des voies métaboliques} : \textbf{+100\%}\\
        $\bullet$ \textbf{Triplets \& Citations d'évidences} : 15 vs \textbf{35}\\
        $\bullet$ Connexions Diabète-Alzheimer : 0 vs 100\% identifiées\\
        $\bullet$ Fichier certifié : \texttt{data/processed/evaluation\_results.json}\\[4pt]
        \begin{tikzpicture}
            \node[badge=EmeraldBorder] {evaluate\_rag.py | Benchmark Suite};
        \end{tikzpicture}
    };

    \node[innerCard=EmeraldBorder, fill=EmeraldPrimary, text width=6.2cm] (llm_box3) at (19.8, 1.8) {
        \textbf{\textcolor{EmeraldText}{\large Interfaces \& Déploiement}}\\[3pt]
        $\bullet$ \textbf{Application Web Streamlit} (\texttt{app.py} \& \texttt{run\_app.bat}) :\\
        \quad - Chatbot GraphRAG avec streaming \& citations\\
        \quad - Visualiseur (36 gènes partagés \& mode de vie)\\
        \quad - Dashboard comparatif Standard RAG vs GraphRAG\\
        $\bullet$ \textbf{CLI Interactif} (\texttt{src/rag/cli.py}) \& Démo (\texttt{demo.py})\\
        $\bullet$ Dépôt GitHub : \texttt{biomedical-graphrag}\\[4pt]
        \begin{tikzpicture}
            \node[badge=EmeraldBorder] {Streamlit (app.py) | CLI | GitHub CI};
        \end{tikzpicture}
    };

    \begin{scope}[on background layer]
        \node[stepCard=EmeraldBorder, fit={(llm_box1) (llm_box2) (llm_box3)}, label={[EmeraldBorder, font=\bfseries\normalsize]above:ÉTAPE 4 \& 5 : RAISONNEMENT CAUSAL MULTI-ÉCHELLES, BENCHMARK \& DÉPLOIEMENT STREAMLIT}] (card_step4) {};
    \end{scope}

    % -------------------------------------------------------------------------
    % FLÈCHES DE FLUX DIRECTIONNELLES (DATAFLOW)
    % -------------------------------------------------------------------------
    % ETL vers Graphe
    \draw[flowArrow=TealBorder] (etl_sub1.east) -- (kg_box.west)
        node[midway, flowLabel=TealBorder] {Triplets CSV};

    % ETL vers Vector Store
    \draw[flowArrow=IndigoBorder] (etl_sub3.east) -- (vec_box.west)
        node[midway, flowLabel=IndigoBorder] {Chunks XML};

    % Graphe vers Hybrid Retriever
    \draw[flowArrow=TealBorder] (kg_box.east) -- (rag_box1.west)
        node[midway, flowLabel=TealBorder] {Chemins $k$-Hop \& 36 Gènes};

    % Vector Store vers Hybrid Retriever
    \draw[flowArrow=IndigoBorder] (vec_box.east) -- (rag_box2.west)
        node[midway, flowLabel=IndigoBorder] {Top-$k$ Chunks (PubMedBERT)};

    % Hybrid Retriever vers Causal Generator
    \draw[flowArrow=AmberBorder] (card_step3.south) -- ++(0,-0.4) -| (llm_box3.north)
        node[pos=0.25, flowLabel=AmberBorder] {Contexte Causal Fusionné};

\end{tikzpicture}

\end{document}
.
% ==============================================================================

\documentclass[border=15pt, 11pt]{standalone}
\usepackage[utf8]{inputenc}
\usepackage[T1]{fontenc}
\usepackage{lmodern}
\usepackage{amsmath, amssymb}
\usepackage{xcolor}
\usepackage{tikz}

% Bibliothèques TikZ nécessaires
\usetikzlibrary{shapes.geometric, arrows.meta, positioning, calc, fit, backgrounds, shadows.blur}

% Palette de couleurs professionnelles (Charte Slate Navy, Teal, Indigo & Emerald)
\definecolor{SlateNavy}{HTML}{0B1120}
\definecolor{SlateCard}{HTML}{0F172A}
\definecolor{SlateSubCard}{HTML}{1E293B}
\definecolor{TealPrimary}{HTML}{042F2E}
\definecolor{TealBorder}{HTML}{2DD4BF}
\definecolor{TealText}{HTML}{5EEAD4}
\definecolor{IndigoPrimary}{HTML}{1E1B4B}
\definecolor{IndigoBorder}{HTML}{818CF8}
\definecolor{IndigoText}{HTML}{A5B4FC}
\definecolor{SkyPrimary}{HTML}{0369A1}
\definecolor{SkyBorder}{HTML}{38BDF8}
\definecolor{AmberPrimary}{HTML}{451A03}
\definecolor{AmberBorder}{HTML}{F59E0B}
\definecolor{AmberText}{HTML}{FDE68A}
\definecolor{EmeraldPrimary}{HTML}{022C22}
\definecolor{EmeraldBorder}{HTML}{10B981}
\definecolor{EmeraldText}{HTML}{6EE7B7}
\definecolor{TextLight}{HTML}{F8FAFC}
\definecolor{TextMuted}{HTML}{94A3B8}
\definecolor{BadgeBg}{HTML}{030712}

\begin{document}

\begin{tikzpicture}[
    font=\sffamily,
    >=Stealth,
    node distance=1.2cm and 1.5cm,
    % Styles des blocs
    titleBox/.style={
        text=TextLight,
        align=center,
        font=\bfseries\Large
    },
    subtitleBox/.style={
        text=SkyBorder,
        align=center,
        font=\bfseries\normalsize
    },
    metaBox/.style={
        text=TextMuted,
        align=center,
        font=\footnotesize\itshape
    },
    stepCard/.style={
        draw=#1,
        line width=1.5pt,
        rounded corners=8pt,
        inner sep=12pt,
        fill=SlateCard,
        blur shadow={shadow blur steps=5, shadow xshift=2pt, shadow yshift=-2pt, shadow opacity=40}
    },
    innerCard/.style={
        draw=#1,
        line width=1pt,
        rounded corners=5pt,
        inner sep=8pt,
        fill=SlateSubCard,
        text=TextLight,
        align=left,
        font=\footnotesize
    },
    badge/.style={
        draw=#1,
        line width=0.8pt,
        rounded corners=3pt,
        fill=BadgeBg,
        text=#1,
        font=\bfseries\scriptsize,
        inner sep=3pt
    },
    flowArrow/.style={
        ->,
        line width=2pt,
        draw=#1,
        shorten >=3pt,
        shorten <=3pt
    },
    flowLabel/.style={
        font=\bfseries\scriptsize,
        fill=SlateNavy,
        inner sep=2pt,
        rounded corners=2pt,
        text=#1
    }
]

    % -------------------------------------------------------------------------
    % TITRE ET EN-TÊTE PRINCIPAL
    % -------------------------------------------------------------------------
    \node[titleBox] (mainTitle) at (10, 15.2) {ARCHITECTURE DU PIPELINE HYBRIDE GRAPHRAG BIOMÉDICAL};
    \node[subtitleBox, below=0.15cm of mainTitle] (subTitle) {
        Hamza Essamoudi $\bullet$ Salma Ouahman $\bullet$ Ghita El Hajjouji \quad$|$\quad Projet PFA --- 5IADA (EMSI)
    };
    \node[metaBox, below=0.12cm of subTitle] (metaLine) {
        Analyse Causale Multi-Échelles : Maladie d'Alzheimer (DOID:10652) $\bullet$ Diabète Type 2 (DOID:9352) $\bullet$ Ontologie Human Lifestyle
    };

    % -------------------------------------------------------------------------
    % ÉTAPE 1 : ACQUISITION & ETL MULTI-ONTOLOGIES
    # -------------------------------------------------------------------------
    \node[innerCard=SkyBorder, text width=6.2cm] (etl_sub1) at (3, 11.2) {
        \textbf{\textcolor{SkyBorder}{\normalsize 1.1 Réseaux Biomédicaux}}\\[3pt]
        $\bullet$ \textbf{PrimeKG-Plus} : Ontologies \& Mappings UMLS/MeSH\\
        $\bullet$ \textbf{Hetionet v1.0} : 2,25M arêtes relationnelles\\
        $\bullet$ Harmonisation : DOID, Entrez, MeSH, Reactome\\[4pt]
        \begin{tikzpicture}
            \node[badge=SkyBorder] {Python 3.12 | Pandas | Git LFS};
        \end{tikzpicture}
    };

    \node[innerCard=TealBorder, text width=6.2cm, below=0.45cm of etl_sub1] (etl_sub2) {
        \textbf{\textcolor{TealBorder}{\normalsize 1.2 Human Lifestyle Ontology (Nouveau)}}\\[3pt]
        $\bullet$ \textbf{15 N\oe uds} hiérarchiques (MeSH \texttt{D008019})\\
        $\bullet$ \textbf{40 Arêtes causales} (\texttt{PROTECTS}, \texttt{AGGRAVATES})\\
        $\bullet$ Facteurs modifiables : Diet, Exercise, Sleep, Stress\\[4pt]
        \begin{tikzpicture}
            \node[badge=TealBorder] {MeSH D008019 | UMLS C0023674 | enrich\_lifestyle.py};
        \end{tikzpicture}
    };

    \node[innerCard=SkyBorder, text width=6.2cm, below=0.45cm of etl_sub2] (etl_sub3) {
        \textbf{\textcolor{SkyBorder}{\normalsize 1.3 Littérature Médicale Certifiée}}\\[3pt]
        $\bullet$ \textbf{NIH MedQuAD} : 47 457 Q/R cliniques révisées\\
        $\bullet$ Dossiers ciblés : NINDS, GHR, NIDDK, SeniorHealth\\
        $\bullet$ Chunking contextualisé : \textbf{719 Chunks} (190 Lifestyle)\\[4pt]
        \begin{tikzpicture}
            \node[badge=SkyBorder] {XML Parsing | Chunking Contextualisé};
        \end{tikzpicture}
    };

    \begin{scope}[on background layer]
        \node[stepCard=SkyBorder, fit={(etl_sub1) (etl_sub2) (etl_sub3)}, label={[SkyBorder, font=\bfseries\normalsize]above:ÉTAPE 1 : ACQUISITION \& ETL MULTI-ONTOLOGIES}] (card_step1) {};
    \end{scope}

    % -------------------------------------------------------------------------
    % ÉTAPE 2A : KNOWLEDGE GRAPH CAUSAL (NEO4J)
    % -------------------------------------------------------------------------
    \node[innerCard=TealBorder, fill=TealPrimary, text width=6.8cm] (kg_box) at (11.5, 11.5) {
        \textbf{\textcolor{TealText}{\large 2A. Knowledge Graph Causal (Neo4j)}}\\[3pt]
        $\bullet$ \textbf{4 518 N\oe uds} uniques (9 labels hiérarchisés)\\
        $\bullet$ \textbf{19 951 Relations} causales \& fonctionnelles (17 types)\\
        $\bullet$ \textbf{Axe Diabète-Alzheimer} : 36 gènes partagés (\textit{Type 3 Diabetes})\\
        $\bullet$ Ingestion batch \texttt{UNWIND} \& Contraintes d'unicité\\
        $\bullet$ Traçabilité $k$-hop : Lifestyle $\to$ Gènes $\to$ Voies $\to$ Maladie\\[4pt]
        \begin{tikzpicture}
            \node[badge=TealBorder] {Neo4j 5.26 | Cypher | APOC | Docker Compose};
        \end{tikzpicture}
    };

    % -------------------------------------------------------------------------
    % ÉTAPE 2B : BASE VECTORIELLE CLINIQUE DENSE
    % -------------------------------------------------------------------------
    \node[innerCard=IndigoBorder, fill=IndigoPrimary, text width=6.8cm, below=0.6cm of kg_box] (vec_box) {
        \textbf{\textcolor{IndigoText}{\large 2B. Base Vectorielle Clinique Dense}}\\[3pt]
        $\bullet$ \textbf{719 Chunks} contextuels MedQuAD indexés\\
        $\bullet$ Modèle biomédical \textbf{PubMedBERT} (768 dimensions)\\
        $\bullet$ Métadonnées cliniques : \texttt{disease}, \texttt{qtype}, \texttt{focus}\\
        $\bullet$ \textbf{Dual Vector Store} :\\
        \quad - \textbf{ChromaDB} : Moteur embarqué local sans dépendance\\
        \quad - \textbf{Qdrant} : Microservice Docker pour la production\\[4pt]
        \begin{tikzpicture}
            \node[badge=IndigoBorder] {PubMedBERT (768d) | ChromaDB | Qdrant Docker};
        \end{tikzpicture}
    };

    \begin{scope}[on background layer]
        \node[stepCard=TealBorder, fit={(kg_box) (vec_box)}, label={[TealBorder, font=\bfseries\normalsize]above:ÉTAPE 2 : STOCKAGE \& INDEXATION DUAL-CORE}] (card_step2) {};
    \end{scope}

    % -------------------------------------------------------------------------
    % ÉTAPE 3 : RETRIEVER HYBRIDE & FUSION DE CONTEXTE
    % -------------------------------------------------------------------------
    \node[innerCard=AmberBorder, fill=AmberPrimary, text width=6.6cm] (rag_box1) at (19.8, 11.5) {
        \textbf{\textcolor{AmberText}{\large 3.1 Retrieval Conjoint}}\\[3pt]
        $\bullet$ \textbf{BiomedicalEntityLinker} : 30+ alias lifestyle \& gènes\\
        $\bullet$ \textbf{Graphe $k$-Hop} : (Lifestyle) $\to$ (BDNF/TNF) $\to$ (Voies)\\
        $\bullet$ \textbf{Axe Diabète-Alzheimer} : 36 gènes communs (IRS1, AKT1)\\
        $\bullet$ \textbf{Recherche Vectorielle} : Top-$k$ chunks certifiés NIH\\
        $\bullet$ Score combiné normalisé de pertinence\\[4pt]
        \begin{tikzpicture}
            \node[badge=AmberBorder] {HybridRetriever | LocalGraphEngine};
        \end{tikzpicture}
    };

    \node[innerCard=AmberBorder, fill=AmberPrimary, text width=6.6cm, below=0.6cm of rag_box1] (rag_box2) {
        \textbf{\textcolor{AmberText}{\large 3.2 Fusion de Contexte}}\\[3pt]
        $\bullet$ Linéarisation ordonnée des triplets \& chaînes causales\\
        $\bullet$ Fusion conjointe : Graphe formel + Extraits certifiés NIH\\
        $\bullet$ Élimination documentée des hallucinations médicales\\
        $\bullet$ Éradication des ruptures déductives et faux liens\\[4pt]
        \begin{tikzpicture}
            \node[badge=AmberBorder] {ContextFusionEngine | Pydantic Schema};
        \end{tikzpicture}
    };

    \begin{scope}[on background layer]
        \node[stepCard=AmberBorder, fit={(rag_box1) (rag_box2)}, label={[AmberBorder, font=\bfseries\normalsize]above:ÉTAPE 3 : MOTEUR HYBRIDE GRAPHRAG}] (card_step3) {};
    \end{scope}

    % -------------------------------------------------------------------------
    % ÉTAPES 4 & 5 : RAISONNEMENT CAUSAL, BENCHMARK & APPLICATION STREAMLIT
    % -------------------------------------------------------------------------
    \node[innerCard=EmeraldBorder, fill=EmeraldPrimary, text width=6.2cm] (llm_box1) at (3.2, 1.8) {
        \textbf{\textcolor{EmeraldText}{\large 4. Moteur de Synthèse Causale}}\\[3pt]
        1. Déclencheurs étiologiques (APP, PSEN1, APOE)\\
        2. Cascades intermédiaires (Résistance insuline, TNF)\\
        3. Impact Lifestyle : Neuroprotection (BDNF)\\
        4. Phénotypes \& déclin cognitif (Démence, mémoire)\\
        5. Repositionnement (Donepezil, Memantine, Diète)\\
        6. Traçabilité absolue : Triplets graphe \& DocIDs NIH\\[4pt]
        \begin{tikzpicture}
            \node[badge=EmeraldBorder] {CausalGenerator | LLM Routing};
        \end{tikzpicture}
    };

    \node[innerCard=EmeraldBorder, fill=EmeraldPrimary, text width=6.6cm] (llm_box2) at (11.5, 1.8) {
        \textbf{\textcolor{EmeraldText}{\large 5. Benchmark d'Évaluation}}\\[3pt]
        $\bullet$ \textbf{Chaînes causales tracées} : 0 (Std) vs \textbf{23} (GraphRAG)\\
        $\bullet$ \textbf{Couverture des voies métaboliques} : \textbf{+100\%}\\
        $\bullet$ \textbf{Triplets \& Citations d'évidences} : 15 vs \textbf{35}\\
        $\bullet$ Connexions Diabète-Alzheimer : 0 vs 100\% identifiées\\
        $\bullet$ Fichier certifié : \texttt{data/processed/evaluation\_results.json}\\[4pt]
        \begin{tikzpicture}
            \node[badge=EmeraldBorder] {evaluate\_rag.py | Benchmark Suite};
        \end{tikzpicture}
    };

    \node[innerCard=EmeraldBorder, fill=EmeraldPrimary, text width=6.2cm] (llm_box3) at (19.8, 1.8) {
        \textbf{\textcolor{EmeraldText}{\large Interfaces \& Déploiement}}\\[3pt]
        $\bullet$ \textbf{Application Web Streamlit} (\texttt{app.py} \& \texttt{run\_app.bat}) :\\
        \quad - Chatbot GraphRAG avec streaming \& citations\\
        \quad - Visualiseur (36 gènes partagés \& mode de vie)\\
        \quad - Dashboard comparatif Standard RAG vs GraphRAG\\
        $\bullet$ \textbf{CLI Interactif} (\texttt{src/rag/cli.py}) \& Démo (\texttt{demo.py})\\
        $\bullet$ Dépôt GitHub : \texttt{biomedical-graphrag}\\[4pt]
        \begin{tikzpicture}
            \node[badge=EmeraldBorder] {Streamlit (app.py) | CLI | GitHub CI};
        \end{tikzpicture}
    };

    \begin{scope}[on background layer]
        \node[stepCard=EmeraldBorder, fit={(llm_box1) (llm_box2) (llm_box3)}, label={[EmeraldBorder, font=\bfseries\normalsize]above:ÉTAPE 4 \& 5 : RAISONNEMENT CAUSAL MULTI-ÉCHELLES, BENCHMARK \& DÉPLOIEMENT STREAMLIT}] (card_step4) {};
    \end{scope}

    % -------------------------------------------------------------------------
    % FLÈCHES DE FLUX DIRECTIONNELLES (DATAFLOW)
    % -------------------------------------------------------------------------
    % ETL vers Graphe
    \draw[flowArrow=TealBorder] (etl_sub1.east) -- (kg_box.west)
        node[midway, flowLabel=TealBorder] {Triplets CSV};

    % ETL vers Vector Store
    \draw[flowArrow=IndigoBorder] (etl_sub3.east) -- (vec_box.west)
        node[midway, flowLabel=IndigoBorder] {Chunks XML};

    % Graphe vers Hybrid Retriever
    \draw[flowArrow=TealBorder] (kg_box.east) -- (rag_box1.west)
        node[midway, flowLabel=TealBorder] {Chemins $k$-Hop \& 36 Gènes};

    % Vector Store vers Hybrid Retriever
    \draw[flowArrow=IndigoBorder] (vec_box.east) -- (rag_box2.west)
        node[midway, flowLabel=IndigoBorder] {Top-$k$ Chunks (PubMedBERT)};

    % Hybrid Retriever vers Causal Generator
    \draw[flowArrow=AmberBorder] (card_step3.south) -- ++(0,-0.4) -| (llm_box3.north)
        node[pos=0.25, flowLabel=AmberBorder] {Contexte Causal Fusionné};

\end{tikzpicture}

\end{document}
